\documentclass[12pt,a4paper]{article}
\usepackage[utf8]{inputenc}
\usepackage[margin=1in]{geometry}
\usepackage{amsmath}
\usepackage{amssymb}
\usepackage{graphicx}
\usepackage{hyperref}
\usepackage{cite}

\title{VMC Automatic Synthesis: Genetic Algorithm for Tendon-Driven Robotic Manipulation}
% titles should be different

\author{}
\date{\today}

\begin{document}
\maketitle


\section*{1. Introduction}
\begin{itemize}
  \item Background and motivation
  \begin{itemize}
    \item Robot manipulation challenges
    \item Virtual model control
    \item motivation for automatic synthesis 
  \end{itemize}
  \item Problem statement
  \begin{itemize}
    \item Project aims
    \item Research questions 
    % can we do it
    % what is the sim to real gap
    % what is the sensitivity 
    \item Scope and limitations
  \end{itemize}
\end{itemize}

\section*{2. Literature Review}
\begin{itemize} 

    % start with what is used to solve the problem
    % design by hand vs ...
    % why using genetic algorithms
    % novelty of how people have addressed this. (designing a mechanism vs trajectory)
  \item Virtual model control principles and applications.
  \item Virtual model control design and optimization methods.
  \item robotic manipulation and grasping strategies.
  \item Genetic algorithms and other optimization techniques in robotics.
  \item simulation-to-real transfer and digital twin.
  \item gap analysis and motivation for our approach.
\end{itemize}

\section*{3. System and Infrastructure Overview}
% motivate why we are providing this information (chapter vs appendix)
\begin{itemize}
  \item Hardware setup and control (Sciurus17)
  \item VMC architecture and control flow.
  \item Simulation environment.
    \begin{itemize}
        \item MuJoCo physics engine.
        \item model construction and simulation tooling (`ModelBuilder`).
        \item contact Modeling (friction, solver, soft materials).
        \item tendon modeling, Jacobian mapping, and controller callbacks.
        \item validation and visualization tools.
    \end{itemize}
  \item Object tracking (Aruco). 
  
\end{itemize}

\section*{4. Automatic Synthesis Methodology}
\begin{itemize}
    \item Problem Formulation
    \item Parameterization of tendon configurations and controllers. (Constrining the search space)
    \item Objective function design and scenario randomization.
    \item Optimization method (Genetic Algorithm).
    \begin{itemize}
        \item GA setup (population, selection, elitism, crossover, mutation).
        \item Enconding of tendon configurations
        \item Improved operators (bounded continuous mutation/crossover, integer creep).
        \item Parallel evaluation and runtime considerations.
    \end{itemize}
    \item runtime considerations and parallel evaluation.
    % this is not essential \item Implementation details Checkpointing, reproducibility, and export artifacts (XML/JSON).
\end{itemize}


\section*{5. Real Robot Experiments}
\begin{itemize}
  \item Initial conditions and scenario
  \item Evaluation metrics and data collection (digital twin, video)
  \item adjustments for sim-to-real transfer (contact calibration, controller scaling, jitter, ramp).
  \item signal processing for controller feedback. (smoothing, delay, noise handling).
\end{itemize}

\section*{6. Results} % Don't separate results and discusstion
\begin{itemize}
  \item Simulation Results
  \begin{itemize}
    \item baseline performance (hand-designed configurations).
    \item optimization progress (fitness curves, diversity).
    \item best configurations and their performance.
  \end{itemize}
%   this should go into implementation deails
  \item Real Robot contoller and object tracking Results
  \item Simulation to real transfer results
  \begin{itemize}
    \item performance of best configurations on real robot.
    \item comparison to simulation results.
    \item qualitative observations (stability, adaptability) and unstable phenomena.
    \item visualization of successful grasps and failure modes.
  \end{itemize}
\end{itemize}

% third section in results which is results after real input.

\section*{7. Discussion}
\begin{itemize}
  \item Analysis of search space and optimization landscape. (fitness function)
  \item GA algorithm performance and convergency analysis.
  \begin{itemize}
    \item effectiveness of GA operators and hyperparameter choices.
    \item genome encoding and its impact on search efficiency.
    \item potential improvements (adaptive mutation, multi-objective optimization).
  \end{itemize}
  \item Analysis of discrepancies between simulation and real-world performance and how to address them.
  \begin{itemize}
    \item Importance of accurate contact modeling and calibration.
    \item joint friction and damping effects
    \item controller power scaling
    \item effect of multi-tendon and high damping configurations and instability.
    \item sensitivity to initial conditions and robustness of controllers.
  \end{itemize}
  \item Generalization and robustness of synthesized controllers.
\end{itemize}

\section*{10. Conclusion}
\begin{itemize}
  \item Summary of technical contributions.
  \item Practical impact and lessons learned. (in discussion section)
\end{itemize}

% future work should be more related to the results rather than just doing more. give a starting point for future work rahther than just trying new stuff.
\section*{11. Future Work}
\begin{itemize}
  \item Multi-objective optimization and Pareto analysis.
  \item Incorporation other optimization techniques (Bayesian optimization, reinforcement learning).
  \item wider search space (moving sites, rotational tendons, etc).
  \item Better contact modeling or use of soft materials in simulation.
  \item making a more robust controller software stack (closed loop control, better signal processing, etc) more stable with higher damping.
  \item Extending to richer grasp tasks.
\end{itemize}



\end{document}


% looked at the problem in two regimes: 
% sticky and non-sticky (more normal) contact.
% maybe can run the stickey situation

% normalisation of the damping
% have discussion on sampling time
% maybe apply sampling delay to the simulation


